% ---
% Subsección 2.1.8 de los Antecedentes (nueva en TT2; antes 2.1.9 en el borrador del 2026-10-14).
% Aplica a tres trabajos las cuatro preguntas de control que el proyecto se aplicó a sí mismo (H-046).
% Cada atribución lleva su respaldo verificable con Code/tools/verificar_citas_fichas.py.
% Las cifras de cada trabajo salen de sus tablas, citadas literalmente en los respaldos.
% ---

\subsubsection{Revisión crítica de los trabajos que reportan ventaja}
\label{sec:revision_critica}

Varios trabajos de aprendizaje automático cuántico aplicado a oncología reportan mejoras frente a clasificadores clásicos. Para valorar esas afirmaciones con el mismo criterio que este trabajo aplica a sus propios resultados, se plantean cuatro preguntas de control:

\begin{enumerate}
  \item \textbf{Paridad}: ¿el clasificador clásico de referencia recibió las mismas variables de entrada, el mismo preprocesamiento y un ajuste de hiperparámetros comparable al del modelo cuántico?
  \item \textbf{Incertidumbre}: ¿la diferencia reportada supera la variabilidad experimental, con desviaciones, intervalos de confianza o pruebas de significación?
  \item \textbf{Partición de prueba}: ¿el conjunto de prueba se usó una sola vez, sin seleccionar modelos ni configuraciones sobre él?
  \item \textbf{Simulabilidad}: ¿el circuito o kernel cuántico evaluado es difícil de simular en una computadora clásica?
\end{enumerate}

La cuarta pregunta tiene un carácter distinto. La dificultad de simulación es una propiedad de una familia de circuitos cuando crece el número de qubits, no de un experimento concreto: con una docena de qubits, el vector de estado tiene unos pocos miles de amplitudes y cualquier circuito se simula de forma exacta. Ningún experimento a esta escala, incluido el presente, puede responderla afirmativamente; se incluye para dejar claro que una ventaja observada a esa escala no se debe a una dificultad de cómputo clásica.

\textbf{Xiang et al. (2024): QCCNN sobre datos tabulares de cáncer de mama}

% respaldo: «In this work, we present a quantum hybrid classical convolutional neural network (QCCNN) based breast cancer diagnosis approach with the goal of utilizing quantum computing’s high-dimensional data processing power and parallelism to increase diagnosis efficiency and accuracy.» (Xiang2024QuantumCNN, p. 1)
Xiang et al. propusieron una red neuronal convolucional híbrida cuántico-clásica (QCCNN) para el diagnóstico del cáncer de mama, con el objetivo de aprovechar el procesamiento en alta dimensión y el paralelismo de la computación cuántica \cite{Xiang2024QuantumCNN}.
\begin{itemize}
  \item \textit{Paridad}: parcial.
% respaldo: «The Kaggle website (a machine learning repository) provided the breast cancer diagnosis GBSG dataset, SEER dataset, and WDBC dataset used in this study.» (Xiang2024QuantumCNN, p. 3)
  El estudio no usa mamografías, sino tres conjuntos tabulares descargados de Kaggle: GBSG, SEER y WDBC \cite{Xiang2024QuantumCNN}.
% respaldo: «a Convolutional Neural Network (CNN) and Logistic Regression model were used for comparative analyses with the QCCNN.» (Xiang2024QuantumCNN, p. 8)
  Los comparadores son una CNN y una regresión logística \cite{Xiang2024QuantumCNN}; no se incluyen otros clasificadores habituales para datos tabulares, como una SVM con kernel RBF o los métodos basados en árboles.
  \item \textit{Incertidumbre}: no.
% respaldo: «The results show that QCCNN performs the best in terms of classification accuracy in all cases, significantly outperforming both the logistic regression model and CNN, confirming its superior performance in breast cancer diagnosis.» (Xiang2024QuantumCNN, p. 11)
  El artículo afirma que QCCNN supera significativamente a ambos modelos en todos los casos \cite{Xiang2024QuantumCNN}, pero sus tablas no reportan desviaciones, intervalos de confianza ni valores $p$.
% respaldo: «Model Training accuracy (%) Validation accuracy (%) QCCNN 74.375 70.498 CNN 59.708 58.944 Logistic regression 61.927 60.583 Table 2.» (Xiang2024QuantumCNN, p. 9)
% respaldo: «Model Training Accuracy (%) Validation Accuracy (%) QCCNN 89.544 90.088 CNN 88.708 89.091 Logistic Regression 89.513 89.125 Table 3.» (Xiang2024QuantumCNN, p. 9)
% respaldo: «Model Training accuracy (%) Validation accuracy (%) QCCNN 97.099 96.155 CNN 92.932 92.822 Logistic regression 91.354 91.184 Table 4.» (Xiang2024QuantumCNN, p. 10)
  La diferencia de exactitud media en validación frente al mejor comparador varía mucho entre conjuntos: 9.9 puntos porcentuales en GBSG (70.5 frente a 60.6~\%), 1.0 en SEER (90.1 frente a 89.1~\%) y 3.3 en WDBC (96.2 frente a 92.8~\%) \cite{Xiang2024QuantumCNN}. Sin una medida de la variabilidad no puede saberse cuáles de esas diferencias superan el ruido muestral.
  \item \textit{Partición de prueba}: no.
% respaldo: «Dataset Sample size Training set Validation set GBSG 686 480 206» (Xiang2024QuantumCNN, p. 4)
% respaldo: «Green indicates the accuracy of the model relative to the training set in each iteration, while purple indicates the accuracy relative to the validation set.» (Xiang2024QuantumCNN, p. 9)
% respaldo: «The highest accuracies of the QCCNN on the training and validation sets are 80.833% and 75.728%, respectively» (Xiang2024QuantumCNN, p. 9)
  Cada conjunto se divide en entrenamiento y validación, sin un conjunto de prueba independiente. Además del promedio, el artículo destaca la exactitud más alta alcanzada en validación a lo largo de las iteraciones \cite{Xiang2024QuantumCNN}; el máximo sobre la misma partición que se reporta es una estimación optimista.
  \item \textit{Simulabilidad}: no.
% respaldo: «we use the PyTorch framework to build and train the classical neural network part and integrate the PennyLane library to implement the quantum convolutional layer.» (Xiang2024QuantumCNN, p. 8)
% respaldo: «In addition, the experiments in this study were conducted on a conventional computer, and although GPUs are able to accelerate some of the computations, their computational speed for quantum circuits is not yet comparable to that of a real quantum computer (QPU), and the use of QPUs will be further explored in future studies to fully evaluate the performance of the QCCNN in practical applications.» (Xiang2024QuantumCNN, p. 11)
  La capa cuántica se simuló con PennyLane sobre PyTorch, en una computadora convencional; los autores dejan el uso de procesadores cuánticos para trabajos futuros \cite{Xiang2024QuantumCNN}.
\end{itemize}

\textbf{El Hafidi et al. (2025): QSVM en cáncer de pulmón}

% respaldo: «QSVM models were implemented using Qiskit and executed on the qasm simulator, employing three distinct quantum feature maps: ZFeatureMap, ZZFeatureMap, and PauliFeatureMap.» (ElHafidi2025, p. 1)
El Hafidi et al. implementaron máquinas de vectores de soporte cuánticas (QSVM) en Qiskit, ejecutadas en el simulador \texttt{qasm\_simulator}, con tres mapas de características: ZFeatureMap, ZZFeatureMap y PauliFeatureMap \cite{ElHafidi2025}.
\begin{itemize}
  \item \textit{Paridad}: sí.
% respaldo: «Using a real-world dataset of 309 patient records with significant class imbalance (39 non-cancer vs. 270 cancer cases), we constructed six balanced subsets for robust evaluation.» (ElHafidi2025, p. 1)
% respaldo: «Each quantum model was evaluated on the same six balanced subsets to ensure a fair comparison with its classical counterpart.» (ElHafidi2025, p. 9)
  Sobre un conjunto clínico de 309 pacientes, con 39 casos sin cáncer y 270 con cáncer, construyeron seis subconjuntos balanceados y evaluaron los modelos cuánticos y clásicos sobre los mismos subconjuntos \cite{ElHafidi2025}.
% respaldo: «The performance of the SVM was evaluated using various configurations of C and γ, optimizing these hyperparameters to achieve the best possible model performance.» (ElHafidi2025, p. 4)
  Los hiperparámetros de la SVM clásica se optimizaron \cite{ElHafidi2025}.
  \item \textit{Incertidumbre}: no.
% respaldo: «for ZZFeatureMap, the results were an accuracy of 0.91, precision of 0.87, recall of 0.96, specificity of 0.85, and F1-score of 0.91; and for PauliFeatureMap, the model achieved a consistent performance of 0.96 across all metrics, including accuracy, precision, recall, specificity, and F1-score.» (ElHafidi2025, p. 9)
% respaldo: «The radial basis function (RBF) kernel provided the best results, with an accuracy of 0.95, precision of 0.98, recall of 0.92, specificity of 0.98, and F1-score of 0.95.» (ElHafidi2025, p. 9)
% respaldo: «For the polynomial kernel, the results improved with an accuracy of 0.94, precision of 0.98, recall of 0.90, specificity of 0.98, and F1-score of 0.93.» (ElHafidi2025, p. 9)
  La QSVM alcanzó una exactitud media de 0.96 con PauliFeatureMap y de 0.91 con ZZFeatureMap; la SVM clásica, 0.95 con kernel RBF y 0.94 con kernel polinomial \cite{ElHafidi2025}. El mejor modelo cuántico queda 0.01 por encima del mejor clásico, y ZZFeatureMap queda por debajo de ambos kernels clásicos. No se reportan desviaciones ni intervalos.
% respaldo: «PauliFeatureMap 1 0.875 0.875 0.875 0.875 0.875 2 0.937 1.000 0.875 1.000 0.933 3 0.937 0.889 1.000 0.875 0.941 4 1.000 1.000 1.000 1.000 1.000» (ElHafidi2025, p. 10)
  Los valores por subconjunto de su tabla de resultados son 0.875, 0.937 y 1.000 \cite{ElHafidi2025}. Son múltiplos de 1/16, lo que es consistente con unas 16 muestras de evaluación por subconjunto: un solo error cambiaría la exactitud en 0.0625, más de seis veces la diferencia entre el mejor modelo cuántico y el mejor clásico. Esta lectura es una inferencia a partir de la tabla; el artículo no declara el tamaño de la partición de evaluación.
  \item \textit{Partición de prueba}: no lo dice. Las métricas se promedian sobre los seis subconjuntos, pero el artículo no describe cómo se separaron entrenamiento y evaluación dentro de cada uno.
  \item \textit{Simulabilidad}: no. Con 15 variables de entrada, los mapas de Qiskit usan un qubit por variable, y un circuito de ese tamaño se simula de forma exacta.
% respaldo: «The PauliFeatureMap generalizes feature encoding further by using a combination of Pauli gates (X, Y , Z), introducing both expressivity and complexity.» (ElHafidi2025, p. 12)
% respaldo: «PauliFeatureMap is the most expressive in theory but suffers from high noise sensitivity, particularly bit flips, due to basis-altering operations.» (ElHafidi2025, p. 12)
  Los autores describen PauliFeatureMap como una combinación de compuertas $X$, $Y$ y $Z$, más expresiva en teoría pero más sensible al ruido de volteo de bit \cite{ElHafidi2025}. En el texto no se indica qué cadenas de Pauli se usaron. En Qiskit, PauliFeatureMap es una familia de circuitos definida por esa lista (sección~\ref{sec:paulifeaturemap}), y con la lista por defecto, $Z$ y $ZZ$, el circuito coincide con el de ZZFeatureMap; sin la lista, el circuito evaluado no queda definido.
\end{itemize}

\textbf{Azevedo et al. (2022): transferencia de aprendizaje clásico-cuántica en mamografía}

% respaldo: «Based on the work of Mari et al., we train a set of hybrid classical-quantum neural networks using transfer learning (TL).» (Azevedo2022QuantumTransfer, p. 1)
% respaldo: «Our task was to solve the problem of classifying full-image mammograms into malignant and benign, provided by BCDR.» (Azevedo2022QuantumTransfer, p. 1)
% respaldo: «BCDR: a breast cancer digital repository» (Azevedo2022QuantumTransfer, p. 14)
Azevedo et al. entrenaron redes híbridas clásico-cuánticas mediante transferencia de aprendizaje (\textit{transfer learning}) para clasificar mamografías completas en benignas y malignas, con imágenes del BCDR (\textit{Breast Cancer Digital Repository}) \cite{Azevedo2022QuantumTransfer}.
\begin{itemize}
  \item \textit{Paridad}: parcial.
% respaldo: «These results were on the same dataset, same train-test split and with a fixed seed.» (Azevedo2022QuantumTransfer, p. 12)
  Todos los modelos se evaluaron sobre el mismo conjunto de datos, con la misma partición y una semilla fija \cite{Azevedo2022QuantumTransfer}.
% respaldo: «A′: a pretrained ResNet18 without the final linear layer, serving as an extractor of 512 features.» (Azevedo2022QuantumTransfer, p. 6)
% respaldo: «Figure 7a shows classical results without TL obtained from a ResNet, meaning that the network was trained from scratch.» (Azevedo2022QuantumTransfer, p. 10)
% respaldo: «For simplicity, we used a simple linear transformation as classifier.» (Azevedo2022QuantumTransfer, p. 10)
  Los modelos con transferencia comparten una ResNet18 preentrenada que extrae 512 variables; el modelo clásico sin transferencia se entrenó desde cero. Con transferencia, el clasificador clásico es una transformación lineal, mientras que en el híbrido se ajustó además la profundidad del circuito \cite{Azevedo2022QuantumTransfer}.
  \item \textit{Incertidumbre}: no.
% respaldo: «Accuracy 0.67 0.84 0.84 0.83 0.83 0.84» (Azevedo2022QuantumTransfer, p. 12)
% respaldo: «AUC 0.70 0.58 0.77 0.77 0.74 0.69» (Azevedo2022QuantumTransfer, p. 12)
  En el conjunto de prueba, el modelo clásico con transferencia y los híbridos de profundidad 1 y 4 alcanzan la misma exactitud, 0.84. El AUC pasa de 0.58 en el clásico con transferencia a 0.77 en el híbrido de profundidad 1; el modelo clásico entrenado desde cero tiene un AUC de 0.70 \cite{Azevedo2022QuantumTransfer}.
% respaldo: «The results for every metric are far more consistent between training and testing and between epochs when using the dressed quantum circuit, so much so that AUC value went from 58 to 77%, leading us to believe that a hybrid-quantum network does present some advantage as AUC measures how well the model separates two classes.» (Azevedo2022QuantumTransfer, p. 12)
  Los autores interpretan esa diferencia de AUC como un indicio de alguna ventaja del circuito híbrido \cite{Azevedo2022QuantumTransfer}, pero no reportan intervalos de confianza ni pruebas estadísticas.
% respaldo: «It is too soon to say whether or not QML can be an advantage. Our results only point to a good generalization of complex data which needs further testing.» (Azevedo2022QuantumTransfer, p. 11)
  En sus conclusiones son prudentes: consideran prematuro decir si el QML puede ser una ventaja \cite{Azevedo2022QuantumTransfer}.
  \item \textit{Partición de prueba}: no.
% respaldo: «Table 4 Performance metrics results for test set in all ResNet models tested.» (Azevedo2022QuantumTransfer, p. 12)
% respaldo: «CQ1 has the higher overall results between training and test, namely precision, specificity and AUC score.» (Azevedo2022QuantumTransfer, p. 10)
% respaldo: «No further testing was needed on depth as the AUC value started to degrade.» (Azevedo2022QuantumTransfer, p. 10)
  La profundidad del circuito se eligió comparando las métricas del conjunto de prueba, y la búsqueda se detuvo cuando el AUC empezó a empeorar \cite{Azevedo2022QuantumTransfer}; el artículo no describe un conjunto de validación separado.
% respaldo: «for 20 repetitions of the testing in the same testset.» (Azevedo2022QuantumTransfer, p. 10)
% respaldo: «In Table 5 we show the metrics results for the best classification given by the quantum device among the 20 repetitions.» (Azevedo2022QuantumTransfer, p. 10)
% respaldo: «Loss Precision Recall F1-score Accuracy Specificity Time (hours) 0.47 0.83 0.80 0.81 0.81 0.84 43.7» (Azevedo2022QuantumTransfer, p. 12)
  En el procesador \texttt{ibm\_lagos}, la prueba se repitió 20 veces sobre el mismo conjunto, y la exactitud reportada, 0.81, es la mejor de esas 20 repeticiones \cite{Azevedo2022QuantumTransfer}.
% respaldo: «We fixed this by oversampling class 1 by adding 2 copies of every sample» (Azevedo2022QuantumTransfer, p. 7)
  Además, el desbalance se corrigió añadiendo copias rotadas de cada mamografía maligna \cite{Azevedo2022QuantumTransfer}, y el artículo no aclara si esas copias se generaron antes o después de separar entrenamiento y prueba.
  \item \textit{Simulabilidad}: no.
% respaldo: «Finally, the 4 output states are measured on the classical register using the Pauli-Z matrix and passed to the linear classical layer, producing 2 output states.» (Azevedo2022QuantumTransfer, p. 6)
% respaldo: «whose variational quantum circuits have a depth of 1, 2, 3, and 4, respectively (using the simulator).» (Azevedo2022QuantumTransfer, p. 10)
  El circuito variacional tiene 4 qubits y entre 1 y 4 capas \cite{Azevedo2022QuantumTransfer}, un tamaño que cualquier computadora simula de forma exacta.
\end{itemize}

Ninguno de los tres trabajos responde afirmativamente a las cuatro preguntas. Esto no invalida sus resultados, pero indica que las mejoras que reportan no bastan para atribuir una ventaja a la componente cuántica.
